\chapter{Zero-free half-planes from Chebyshev bounds on windows}\label{chapter-windowed-landau}

For a primitive Dirichlet character $\psi$, this chapter reduces the
statement that $L(s,\psi)\neq0$ for $\operatorname{Re}s>1/2$ to a one-sided
bound $G_\psi(x)\leq x^{1/2+O(\varepsilon)}$, where $G_\psi$ is a positive
combination of the Chebyshev sums $\sum_{n\leq x}\Lambda(n)\chi(n)(x/n-1)$
over the cyclic group generated by $\psi$, each corrected by its main term.
The bound is needed only on windows $[Y,Y^{\lambda}]$ of scales for
arbitrarily large $Y$; the later chapters supply it with $\lambda=3/2$. The
deduction is a quantitative form of Landau's theorem on Laplace transforms
of functions of one sign. The positivity it uses is that a real zero gives
the Laplace transform of $G_\psi(e^u)$ a pole with positive residue.

\section{The weighted Chebyshev sum}\label{section-wl-weighted-sum}

\begin{definition}\label{wl-corrected-sums}
Let $N\geq1$, let $\chi$ be a Dirichlet character modulo $N$
(Definition~\ref{mt-dirichlet-character}), and let $L(s,\chi)$ be its
$L$-function (Definition~\ref{mt-l-function}). Write $1$ for the trivial
character modulo $1$, so that $L(s,1)=\zeta(s)$, and call $\chi$
\emph{nontrivial} if $\chi\neq\mathbf 1_N$. Let $\Lambda$ be the von Mangoldt
function and $\gamma$ Euler's constant. For real $y>0$ put
\[
 G(y)=\sum_{1\leq n\leq y}\Lambda(n)\Bigl(\frac yn-1\Bigr)-y\log y+(1+\gamma)y,
\]
and, if $\chi$ is nontrivial,
\[
 g_\chi(y)=\sum_{1\leq n\leq y}\Lambda(n)\chi(n)\Bigl(\frac yn-1\Bigr)
 +y\,\frac{L'}{L}(1,\chi).
\]
For the trivial character modulo $1$ put $g_1=G$. These are the
\emph{corrected Chebyshev sums}.
\end{definition}
\leanref{Formalism/StatementOC.lean}{Formalism.Gz}
\leanref{Formalism/SharpRestriction/RamificationFrobenius/TotientGaloisClasses/Analytic.lean}{Formalism.HC.An.G}
\leanref{Formalism/SharpRestriction/RamificationFrobenius/TotientGaloisClasses/Analytic.lean}{Formalism.HC.An.gChi}
\leanref{Development/AutomorphicArithmetic/DedekindLFunctions/CharFamily.lean}{Deduction.Gchi}

The sum $\sum_{n\leq y}\Lambda(n)\chi(n)(y/n-1)$ is the inverse Mellin
transform of $-\frac{L'}{L}(s,\chi)/(s(s-1))$. The residue of
$-\frac{L'}{L}(s,\chi)y^s/(s(s-1))$ at $s=1$ is $-y\frac{L'}{L}(1,\chi)$ for
nontrivial $\chi$ and $y\log y-(1+\gamma)y$ for $\zeta$, and the corrections
subtract it. Lemma~\ref{wl-laplace-character} computes the transforms of the
corrected sums.

\begin{definition}\label{wl-character-family}
A \emph{character family} $F=(\chi_i,a_i)_{i\in I}$ consists of a finite set
$I$ and, for each $i\in I$, an integer $N_i\geq1$, a Dirichlet character
$\chi_i$ modulo $N_i$ (Definition~\ref{mt-dirichlet-character}) and a real
number $a_i>0$, such that for each $i$ either $\chi_i$ is nontrivial or
$N_i=1$. Its \emph{weighted Chebyshev sum}
and its \emph{zero set} are
\[
 G_F(y)=\sum_{i\in I}a_i\operatorname{Re}g_{\chi_i}(y)\quad(y>0),\qquad
 Z_F=\{\rho\in\mathbf C:\ \operatorname{Re}\rho>0,\ L(\rho,\chi_i)=0
 \text{ or }L(\bar\rho,\chi_i)=0\text{ for some }i\in I\},
\]
with $g_{\chi_i}$ as in Definition~\ref{wl-corrected-sums}.
\end{definition}
\leanref{Development/AutomorphicArithmetic/DedekindLFunctions/CharFamily.lean}{Deduction.CharFamily}
\leanref{Development/AutomorphicArithmetic/DedekindLFunctions/CharFamily.lean}{Deduction.Gm}
\leanref{Development/AutomorphicArithmetic/DedekindLFunctions/LandauZeroFreeCriterion.lean}{Deduction.zeroSet}

\begin{definition}\label{wl-level}
A primitive Dirichlet character (Definition~\ref{mt-dirichlet-character}) is
regarded as a character modulo its conductor; the trivial one is the character
modulo $1$. A \emph{level} $K$
consists of:
\begin{enumerate}
\item a finite set of primitive Dirichlet characters, also written $K$, that
  contains the trivial character and is closed under $\chi\mapsto\bar\chi$;
\item either no element or one element $\psi_K\in K$, the \emph{exceptional
  character} of $K$, which is odd and takes real values;
\item an integer $o\geq0$.
\end{enumerate}
\end{definition}
\leanref{Formalism/SharpRestriction/RamificationFrobenius/TotientGaloisClasses/PrimePairSizeCondition.lean}{Formalism.HC.PChar}
\leanref{Formalism/SharpRestriction/RamificationFrobenius/TotientGaloisClasses/PrimePairSizeCondition.lean}{Formalism.HC.KSet}
\leanref{Formalism/SharpRestriction/RamificationFrobenius/TotientGaloisClasses/ExplicitFormulaErrorOfHostConditions.lean}{Formalism.HC.Level}

\begin{definition}\label{wl-weighted-sum}
Let $K$ be a level (Definition~\ref{wl-level}) with integer $o$. For
$\chi\in K$ put $w_\chi=1$ if $\chi$ is the exceptional character of $K$, and
$w_\chi=(3o+1)^{-1}$ otherwise. The \emph{weighted Chebyshev sum} of $K$ is
\[
 G_K(y)=G(y)+\sum_{\chi\in K,\ \chi\neq1}w_\chi\operatorname{Re}g_\chi(y)
 \qquad(y>0),
\]
with $G$ and $g_\chi$ as in Definition~\ref{wl-corrected-sums}.
\end{definition}
\leanref{Formalism/SharpRestriction/RamificationFrobenius/TotientGaloisClasses/ExplicitFormulaErrorOfHostConditions.lean}{Formalism.HC.Level.wt}
\leanref{Formalism/SharpRestriction/RamificationFrobenius/TotientGaloisClasses/ExplicitFormulaErrorOfHostConditions.lean}{Formalism.HC.Level.Gm}

\begin{definition}\label{wl-cyclic-level}
Let $\psi$ be a primitive Dirichlet character modulo $N$, and let
$\langle\psi\rangle$ be the cyclic group of Dirichlet characters modulo $N$
generated by $\psi$, of order $o_\psi$. The \emph{cyclic level} $K_\psi$ of
$\psi$ is the level whose set consists of the primitive characters inducing
the elements of $\langle\psi\rangle$, with no exceptional character and with
integer $o=o_\psi$. It contains $\psi$. Its weighted Chebyshev sum
(Definition~\ref{wl-weighted-sum}) is written $G_\psi=G_{K_\psi}$; thus, for
$y>0$,
\[
 G_\psi(y)=\sum_{1\leq n\leq y}\Lambda(n)\Bigl(\frac yn-1\Bigr)-y\log y+(1+\gamma)y
 +\frac{1}{3o_\psi+1}\sum_{\chi\in K_\psi,\ \chi\neq1}
 \Bigl(\sum_{1\leq n\leq y}\Lambda(n)\operatorname{Re}\chi(n)\Bigl(\frac yn-1\Bigr)
 +y\operatorname{Re}\frac{L'}{L}(1,\chi)\Bigr).
\]
\end{definition}
\leanref{Formalism/SharpRestriction/RamificationFrobenius/TotientGaloisClasses/ExplicitFormulaErrorOfHostConditions.lean}{Formalism.HC.Level.ofSub}
\leanref{Formalism/SharpRestriction/RamificationFrobenius/TotientGaloisClasses/ExplicitFormulaErrorOfHostConditions.lean}{Formalism.HC.Level.cyc}
\leanref{Formalism/SharpRestriction/RamificationFrobenius/TotientGaloisClasses/ExplicitFormulaErrorOfHostConditions.lean}{Formalism.HC.Level.mem_cyc}

The set $K_\psi$ is closed under conjugation because the primitive character
inducing $\bar\chi$ is the conjugate of the one inducing $\chi$. The
characters of $K_\psi$ have conductors dividing $N$, and $K_\psi$ has $o_\psi$
elements, since a character modulo $N$ is determined by the primitive
character inducing it. For $\psi$ trivial, $K_\psi=\{1\}$ and $G_\psi=G$.

\section{Laplace transforms of corrected Chebyshev sums}\label{section-wl-laplace}

\begin{lemma}\label{wl-laplace-character}
Let $N\geq1$ and let $\chi$ be a Dirichlet character modulo $N$ that is
nontrivial, or the trivial character modulo $1$. Let $f(s)=L(s,\chi)$ if
$\chi$ is nontrivial and $f(s)=(s-1)\zeta(s)$ if $\chi=1$; then $f$ is entire
and $f(1)\neq0$. For $\rho\in\mathbf C$ let $m(\rho)\geq0$ be the order of
vanishing of $f$ at $\rho$. For $s\neq0,1$ with $f(s)\neq0$ put
\[
 \Phi_\chi(s)=-\frac{\frac{L'}{L}(s,\chi)-s\frac{L'}{L}(1,\chi)}{s(s-1)}
 \quad(\chi\neq1),\qquad
 \Phi_1(s)=-\frac{\frac{\zeta'}{\zeta}(s)+\frac1{s-1}+1-(1+\gamma)s}{s(s-1)}.
\]
Then:
\begin{enumerate}
\item for $\operatorname{Re}s>1$,
  \[
   \Phi_\chi(s)=\int_1^\infty g_\chi(y)y^{-s-1}\,dy
   =\int_0^\infty g_\chi(e^u)e^{-su}\,du,
  \]
  with $g_\chi$ as in Definition~\ref{wl-corrected-sums};
\item $\Phi_\chi$ extends to a meromorphic function on $\mathbf C$ that is
  holomorphic at every $\rho\neq0$ with $m(\rho)=0$, and for every $\rho\neq0$
  the function
  \[
   \Phi_\chi(s)-\frac{-m(\rho)}{\rho(\rho-1)}\cdot\frac1{s-\rho}
  \]
  is holomorphic near $\rho$.
\end{enumerate}
\end{lemma}
\begin{proof}
Let $\operatorname{Re}s>1$. For $n\geq1$,
\[
 \int_n^\infty\Bigl(\frac yn-1\Bigr)y^{-s-1}\,dy=\frac{n^{-s}}{s-1}-\frac{n^{-s}}{s}
 =\frac{n^{-s}}{s(s-1)}.
\]
The series $-\frac{L'}{L}(s,\chi)=\sum_n\Lambda(n)\chi(n)n^{-s}$ converges
absolutely, so summing over $n$ gives
$\int_1^\infty\sum_{n\leq y}\Lambda(n)\chi(n)(y/n-1)y^{-s-1}\,dy
=-\frac{L'}{L}(s,\chi)/(s(s-1))$. Also $\int_1^\infty y^{-s}\,dy=1/(s-1)$ and
$\int_1^\infty y^{-s}\log y\,dy=1/(s-1)^2$. For $\chi\neq1$ these give the
first formula. For $\chi=1$ they give
$-\frac{\zeta'}{\zeta}(s)/(s(s-1))-1/(s-1)^2+(1+\gamma)/(s-1)$, which equals
$\Phi_1(s)$ because $1/(s(s-1)^2)+1/(s(s-1))=1/(s-1)^2$. The substitution
$y=e^u$ gives the Laplace transform.

For (2), in both cases
\[
 \Phi_\chi(s)=-\frac{\frac{f'}{f}(s)-s\frac{f'}{f}(1)+k(1-s)}{s(s-1)},
\]
with $k=0$ for $\chi\neq1$ and $k=1$ for $\chi=1$. For the latter,
$\frac{f'}{f}(s)=\frac1{s-1}+\frac{\zeta'}{\zeta}(s)$, and $f(1)=1$,
$f'(1)=\gamma$ from $\zeta(s)=\frac1{s-1}+\gamma+O(s-1)$. The numerator is
holomorphic where $f\neq0$ and vanishes at $s=1$, so $\Phi_\chi$ is
holomorphic at every $\rho\neq0$ with $f(\rho)\neq0$, including $\rho=1$. At
a zero $\rho\neq0$ of $f$ of order $m$, the function
$\frac{f'}{f}(s)-\frac m{s-\rho}$ is holomorphic near $\rho$ and $\rho\neq1$,
so $\Phi_\chi$ has a simple pole at $\rho$ with residue $-m/(\rho(\rho-1))$.
\end{proof}
\leanref{Development/AutomorphicArithmetic/DedekindLFunctions/LandauTheorems/mellinGchiOfNeOne.lean}{Deduction.mellin_Gchi_of_ne_one}
\leanref{Development/AutomorphicArithmetic/DedekindLFunctions/LandauTheorems/mellinGchiModOne.lean}{Deduction.mellin_Gchi_modOne}
\leanref{Development/AutomorphicArithmetic/DedekindLFunctions/LandauTheorems/poleData.lean}{Deduction.poleData_abstract}
\leanref{Development/AutomorphicArithmetic/DedekindLFunctions/LandauTheorems/poleData.lean}{Deduction.poleData_char}

\begin{definition}\label{wl-pole-data}
Let $g:\mathbf R\to\mathbf R$ and $\Phi:\mathbf C\to\mathbf C$ be functions
and let $Z\subset\mathbf C$. The triple $(g,\Phi,Z)$ is \emph{pole data} if:
\begin{enumerate}
\item $g$ is measurable, and there is $C_0$ with $|g(u)|\leq C_0e^{2u}$ for
  all $u\geq0$;
\item $\Phi(s)=\int_0^\infty g(u)e^{-su}\,du$ for $\operatorname{Re}s>2$;
\item $\operatorname{Re}\rho>0$ for every $\rho\in Z$, and every
  $z\in\mathbf C$ with $\operatorname{Re}z>0$ has a punctured neighbourhood
  disjoint from $Z$;
\item $\Phi$ is holomorphic on $\{s:\operatorname{Re}s>0\}\setminus Z$;
\item for every $\rho\in Z$ there are $c\in\mathbf C$ with $c\neq0$ and a
  function $h$ holomorphic near $\rho$ such that $\Phi(s)=c/(s-\rho)+h(s)$ for
  $s\neq\rho$ near $\rho$; if $\rho$ is real, then $c$ is real and $c>0$.
\end{enumerate}
\end{definition}
\leanref{Development/AutomorphicArithmetic/DedekindLFunctions/LandauZeroFreeCriterion.lean}{Deduction.PoleData}

\begin{proposition}\label{wl-pole-data-family}
Let $F=(\chi_i,a_i)_{i\in I}$ be a character family with weighted Chebyshev
sum $G_F$ and zero set $Z_F$ (Definition~\ref{wl-character-family}), and put
$g_F(u)=G_F(e^u)$. Then there is a function $\Phi_F:\mathbf C\to\mathbf C$
such that $(g_F,\Phi_F,Z_F)$ is pole data (Definition~\ref{wl-pole-data}).
\end{proposition}
\begin{proof}
For each $i$, the characters $\chi_i$ and $\bar\chi_i$ satisfy the hypothesis
of Lemma~\ref{wl-laplace-character}. Put
\[
 \Phi_F=\sum_{i\in I}\frac{a_i}2\bigl(\Phi_{\chi_i}+\Phi_{\bar\chi_i}\bigr)
\]
away from the poles of the summands, and give $\Phi_F$ arbitrary values at
those poles.

The function $g_F$ is measurable. Since $\Lambda(n)\leq n$, each $g_{\chi_i}$
satisfies $|g_{\chi_i}(y)|\leq C_iy^2$ for $y\geq1$, so
$|g_F(u)|\leq C_0e^{2u}$ for $u\geq0$. From
$L(\bar s,\bar\chi)=\overline{L(s,\chi)}$ we get
$\frac{L'}{L}(1,\bar\chi)=\overline{\frac{L'}{L}(1,\chi)}$, hence
$g_{\bar\chi}(y)=\overline{g_\chi(y)}$ for real $y$, and
$\operatorname{Re}g_\chi=\frac12(g_\chi+g_{\bar\chi})$. So
Lemma~\ref{wl-laplace-character}(1) gives the Laplace transform for
$\operatorname{Re}s>2$, the integrals converging absolutely by the growth
bound.

The same identity shows that $L(s,\bar\chi_i)=0$ if and only if
$L(\bar s,\chi_i)=0$. So at a point $s\notin Z_F$ with $\operatorname{Re}s>0$
no $f$ of Lemma~\ref{wl-laplace-character} for $\chi_i$ or $\bar\chi_i$
vanishes, and $\Phi_F$ is holomorphic there. Each of these $f$ is a nonzero
entire function, so its zeros are isolated, and so is $Z_F$ in the sense of
Definition~\ref{wl-pole-data}(3).

Let $\rho\in Z_F$. Since $L(s,\chi)\neq0$ for $\operatorname{Re}s\geq1$ and
every $\chi$ in the lemma \cite[Chapters 13 and 14]{Dav00}, we have
$\operatorname{Re}\rho<1$; in particular $\rho\neq1$, and $\rho\neq0$. By
Lemma~\ref{wl-laplace-character}(2), the function $\Phi_F(s)-c/(s-\rho)$ is
holomorphic near $\rho$, where
\[
 c=-\frac{S}{\rho(\rho-1)},\qquad
 S=\sum_{i\in I}\frac{a_i}2\bigl(m_{\chi_i}(\rho)+m_{\bar\chi_i}(\rho)\bigr),
\]
and $m_\chi(\rho)$ is the order of vanishing of $L(s,\chi)$ at $\rho$. As
$\rho\in Z_F$, some $m_{\chi_i}(\rho)$ or $m_{\bar\chi_i}(\rho)$ is positive,
and all $a_i>0$, so $S>0$ and $c\neq0$. If $\rho$ is real, then
$0<\rho<1$, so $\rho(\rho-1)<0$ and $c=S/(\rho(1-\rho))>0$.
\end{proof}
\leanref{Development/AutomorphicArithmetic/DedekindLFunctions/LandauTheorems/poleData.lean}{Deduction.poleData}
\leanref{Development/AutomorphicArithmetic/DedekindLFunctions/LandauTheorems/GchiConj.lean}{Deduction.Gchi_conj}
\leanref{Development/AutomorphicArithmetic/DedekindLFunctions/LandauTheorems/LFunctionConj.lean}{Deduction.LFunction_conj}
\leanref{Development/AutomorphicArithmetic/DedekindLFunctions/LandauTheorems/GchiGrowth.lean}{Deduction.Gchi_growth}
\leanref{Development/AutomorphicArithmetic/DedekindLFunctions/LandauTheorems/CharFamilyLFunctionNeZeroOfOneLeRe.lean}{Deduction.charFamily_LFunction_ne_zero_of_one_le_re}

\section{A quantitative Landau theorem}\label{section-wl-landau}

Landau's theorem \cite{Lan05} says that the Laplace transform of a function
that is eventually of one sign has a singularity at the real point of its
line of convergence \cite[Chapter V]{Ing32}. Here $g$ is only bounded above,
and only on windows $[U,(1+\delta)U]$ for a sequence of $U$ tending to
infinity. The proof works with the Taylor coefficients of the Laplace
transform at a point $\sigma+i\gamma$ to the right of all singularities.

\begin{definition}\label{wl-window-hypothesis}
Let $g:\mathbf R\to\mathbf R$ be a function, let $\theta\in\mathbf R$ and let
$\delta>0$. The function $g$ satisfies the \emph{window hypothesis} with
exponent $\theta$ and width $\delta$ if for every $\eta>0$ and every
$U_0\in\mathbf R$ there are $U\geq U_0$ and a function $R$ on
$[U,(1+\delta)U]$ such that $R(u)e^{-\theta u}$ is integrable on
$[U,(1+\delta)U]$, $0\leq R(u)$ and $g(u)\leq R(u)$ for every
$u\in[U,(1+\delta)U]$, and
\[
 \int_U^{(1+\delta)U}R(u)e^{-\theta u}\,du\leq e^{\eta U}.
\]
\end{definition}
\leanref{Development/AutomorphicArithmetic/DedekindLFunctions/LandauZeroFreeCriterion.lean}{Deduction.WindowHypU}

\begin{lemma}\label{wl-kernel-tails}
Let $\delta>0$. There are $\sigma>2$ and $C_1\in\mathbf R$ such that for every
$U\geq1$ there is an integer $k\geq U$ with
\[
 \int_{(0,\infty)\setminus[U,(1+\delta)U]}\frac{u^k}{k!}e^{-(\sigma-2)u}\,du
 \leq C_1\sigma^{-k}.
\]
\end{lemma}
\begin{proof}
Put $l=2+6/\delta$, so that $l>2$ and $l\delta=6+2\delta$. Put
$M=l(1+\delta)$, $\sigma=M(2+l)$ and $b=\sigma-2$. Then
$b-l=(M-1)(2+l)>0$, and $l\delta=6+2\delta$ gives
\[
 \Bigl(1+\frac\delta2\Bigr)(2+l)=M-1,\qquad
 \Bigl(1+\frac\delta2\Bigr)(l-2)=l+1.
\]
For $a>0$ we have $\int_0^\infty u^ke^{-au}/k!\,du=a^{-k-1}$. Since
$\sigma/a\leq e^{(\sigma-a)/a}$, for every real $x$,
\[
 e^xa^{-k-1}\leq e^{x+k(\sigma-a)/a}a^{-1}\sigma^{-k}.
\]
Let $U\geq1$, put $V=(1+\delta)U$ and let $k=\lceil(1+\delta/2)\sigma U\rceil$,
so $k\geq U$. For $u\notin[U,V]$ one of $l(u-V)$ and $l(U-u)$ is nonnegative,
so
\[
 e^{-bu}\leq e^{-lV}e^{-(b-l)u}+e^{lU}e^{-(b+l)u}.
\]
Multiplying by $u^k/k!$ and integrating over $u>0$ bounds the integral in
question by $e^{-lV}(b-l)^{-k-1}+e^{lU}(b+l)^{-k-1}$.

For the first term take $a=b-l$, so $(\sigma-a)/a=1/(M-1)$. As
$(1+\delta/2)\sigma=M(M-1)$, we have $k\leq M(M-1)U+1$, and $M-1\geq1$ gives
$k/(M-1)\leq MU+1=lV+1$. So the first term is at most $e\,(b-l)^{-1}\sigma^{-k}$.
For the second take $a=b+l$, so $(\sigma-a)/a=-(l-2)/(b+l)$. We have
$k(l-2)\geq(1+\delta/2)(l-2)\sigma U=(l+1)\sigma U$, and
$(l+1)\sigma\geq l(b+l)$ because $\sigma=M(2+l)\geq l(l-2)$. So the exponent
$lU-k(l-2)/(b+l)$ is at most $0$, and the second term is at most
$(b+l)^{-1}\sigma^{-k}$. Hence $C_1=e/(b-l)+1/(b+l)$ will do.
\end{proof}
\leanref{Development/AutomorphicArithmetic/DedekindLFunctions/LandauTheorems/landauTails.lean}{Deduction.landau_tails}
\leanref{Development/AutomorphicArithmetic/DedekindLFunctions/LandauTheorems/kernelIntegral.lean}{Deduction.kernel_integral}

\begin{lemma}\label{wl-coefficient-bound}
Let $(g,\Phi,Z)$ be pole data (Definition~\ref{wl-pole-data}). Let
$0<\theta<1$ and $\delta>0$, and suppose that $g$ satisfies the window
hypothesis with exponent $\theta$ and width $\delta$
(Definition~\ref{wl-window-hypothesis}). For $\operatorname{Re}s>2$ and
integers $k\geq0$ put $a_k(s)=(-1)^k\Phi^{(k)}(s)/k!$. Then there are
$\sigma>2$ and $C\in\mathbf R$ such that for every $\eta>0$ and every
$U_0\in\mathbf R$ there are $U\geq U_0$ and an integer $k\geq U$ with
\[
 |a_k(\sigma+i\gamma)|+\operatorname{Re}a_k(\sigma)\leq Ce^{\eta U}(\sigma-\theta)^{-k}
 \qquad\text{for every }\gamma\in\mathbf R.
\]
\end{lemma}
\begin{proof}
Let $C_0\geq0$ with $|g(u)|\leq C_0e^{2u}$ for $u\geq0$, let $\sigma>2$ and
$C_1$ be given by Lemma~\ref{wl-kernel-tails} for $\delta$, and put
$C=2+2C_0|C_1|$. For $\operatorname{Re}s>2$, differentiation under the
integral sign, justified by the growth bound, gives
\[
 a_k(s)=\int_0^\infty g(u)\frac{u^k}{k!}e^{-su}\,du.
\]

Let $\eta>0$ and $U_0$ be given. The window hypothesis gives
$U\geq\max(U_0,1)$ and a function $R$, and Lemma~\ref{wl-kernel-tails} gives
an integer $k\geq U$. Put $V=(1+\delta)U$ and
$w(u)=u^ke^{-\sigma u}/k!\geq0$. For real $\gamma$ we have
$|a_k(\sigma+i\gamma)|\leq\int_0^\infty|g|w$, and $a_k(\sigma)=\int_0^\infty gw$
is real, so
\[
 |a_k(\sigma+i\gamma)|+\operatorname{Re}a_k(\sigma)
 \leq\int_0^\infty\bigl(|g(u)|+g(u)\bigr)w(u)\,du.
\]
For $u\in[U,V]$, $|g(u)|+g(u)\leq2R(u)$ because $R(u)\geq\max(g(u),0)$, and
$w(u)\leq(\sigma-\theta)^{-k}e^{-\theta u}$ because $t^ke^{-t}/k!\leq1$ for
$t=(\sigma-\theta)u\geq0$. So the integral over $[U,V]$ is at most
$2(\sigma-\theta)^{-k}\int_U^VR(u)e^{-\theta u}\,du\leq2(\sigma-\theta)^{-k}e^{\eta U}$.
For $u>0$ outside $[U,V]$, $|g(u)|+g(u)\leq2C_0e^{2u}$, so the integral there
is at most
\[
 2C_0\int_{(0,\infty)\setminus[U,V]}\frac{u^k}{k!}e^{-(\sigma-2)u}\,du
 \leq2C_0C_1\sigma^{-k}\leq2C_0|C_1|e^{\eta U}(\sigma-\theta)^{-k},
\]
since $0<\sigma-\theta<\sigma$ and $e^{\eta U}\geq1$. Adding the two bounds
gives the claim.
\end{proof}
\leanref{Development/AutomorphicArithmetic/DedekindLFunctions/LandauTheorems/landauWindow.lean}{Deduction.landau_window_key}
\leanref{Development/AutomorphicArithmetic/DedekindLFunctions/LandauTheorems/landauWindowIneq.lean}{Deduction.landau_window_ineq}
\leanref{Development/AutomorphicArithmetic/DedekindLFunctions/LandauTheorems/LandauCoeff.lean}{Deduction.landau_coeff}
\leanref{Development/AutomorphicArithmetic/DedekindLFunctions/LandauTheorems/LandauKernelMax.lean}{Deduction.landau_kernel_max}

\begin{lemma}\label{wl-nearest-pole}
Let $Z\subset\mathbf C$ be a set such that every $z$ with
$\operatorname{Re}z>0$ has a punctured neighbourhood disjoint from $Z$. Let
$0<\theta<\sigma$, suppose that $\operatorname{Re}\rho<\sigma$ for every
$\rho\in Z$, and suppose that some $\rho_0\in Z$ has
$\operatorname{Re}\rho_0>\theta$. Write $B(z,r)=\{w\in\mathbf C:|w-z|<r\}$.
Then one of the following holds:
\begin{enumerate}
\item there are $D\leq\sigma-\theta$ with $Z\cap B(\sigma,D)=\emptyset$, a
  real number $\gamma$, a point $\rho\in Z$ and $r\leq D$ with
  $|\sigma+i\gamma-\rho|<r$ and $Z\cap B(\sigma+i\gamma,r)=\{\rho\}$;
\item there are a real number $\beta>\theta$ with $\beta\in Z$ and $r$ with
  $\sigma-\beta<r\leq\sigma-\theta$ and $Z\cap B(\sigma,r)=\{\beta\}$.
\end{enumerate}
\end{lemma}
\begin{proof}
A closed disc about a point $z$ with radius less than $\operatorname{Re}z$ is
a compact subset of the half-plane $\operatorname{Re}>0$, so it meets $Z$ in
a finite set. Let $F=Z\cap\bar B(\sigma,\sigma-\theta)$, which is finite.

Suppose first that there are $D\leq\sigma-\theta$ with
$Z\cap B(\sigma,D)=\emptyset$ and $\rho_1\in Z$ with
$\sigma-\operatorname{Re}\rho_1<D$. Choose $\varepsilon>0$ with
$2\varepsilon\leq D-(\sigma-\operatorname{Re}\rho_1)$ and
$2\varepsilon\leq\sigma-D$, put $c_0=\sigma+i\operatorname{Im}\rho_1$, and let
$F'=Z\cap\bar B(c_0,D+\varepsilon)$, which is finite since
$D+\varepsilon<\sigma$. For distinct $a,b\in F'$ at most one real $\gamma$
has $|\sigma+i\gamma-a|=|\sigma+i\gamma-b|$: if
$\operatorname{Im}a\neq\operatorname{Im}b$ the equation is linear in $\gamma$
with nonzero coefficient, and if $\operatorname{Im}a=\operatorname{Im}b$ it
forces $\operatorname{Re}a=\operatorname{Re}b$, since both are less than
$\sigma$. So there is $\gamma$ with $|\gamma-\operatorname{Im}\rho_1|<\varepsilon$
such that $z=\sigma+i\gamma$ has distinct distances to distinct points of
$F'$. Then $|z-\rho_1|<\varepsilon+\sigma-\operatorname{Re}\rho_1<D$. Let
$\rho\in F'$ be the point nearest to $z$; it is unique, and
$|z-\rho|\leq|z-\rho_1|<D$. Choose $r\leq D$ with $|z-\rho|<r$ and $r$ at most
the distance from $z$ to every other point of $F'$. A point $w\in Z$ with
$|w-z|<r$ has $|w-c_0|<r+\varepsilon\leq D+\varepsilon$, so $w\in F'$ and
$w=\rho$. This is (1).

If $F$ is empty, the hypothesis of the previous paragraph holds with
$D=\sigma-\theta$ and $\rho_1=\rho_0$. Otherwise let $d$ be the least
distance from $\sigma$ to a point of $F$, so that $Z\cap B(\sigma,d)=\emptyset$.
If some $\rho_1\in Z$ has $\sigma-\operatorname{Re}\rho_1<d$, the hypothesis
holds with $D=d$. In the remaining case $\sigma-\operatorname{Re}\rho\geq d$
for every $\rho\in Z$. A point $\rho\in F$ at distance $d$ from $\sigma$ then
satisfies $|\sigma-\rho|=d\leq\sigma-\operatorname{Re}\rho$, which forces
$\rho=\sigma-d$. So $\beta=\sigma-d$ lies in $Z$ and is the only point of $F$
at distance $d$. If $d=\sigma-\theta$, then
$\sigma-\operatorname{Re}\rho_0<d$, which is excluded; so $d<\sigma-\theta$
and $\beta>\theta$. Choose $r$ with $d<r\leq\sigma-\theta$ and $r$ at most the
distance from $\sigma$ to every point of $F$ other than $\beta$. Then
$Z\cap B(\sigma,r)=\{\beta\}$, which is (2).
\end{proof}
\leanref{Development/AutomorphicArithmetic/DedekindLFunctions/LandauTheorems/LandauDichotomy.lean}{Deduction.landau_dichotomy}
\leanref{Development/AutomorphicArithmetic/DedekindLFunctions/LandauTheorems/LandauNearest.lean}{Deduction.landau_nearest}
\leanref{Development/AutomorphicArithmetic/DedekindLFunctions/LandauTheorems/LandauFinite.lean}{Deduction.landau_finite}

\begin{lemma}\label{wl-pole-coefficients}
Let $z\in\mathbf C$, let $0<r'<r$, and write $B(z,r)=\{w:|w-z|<r\}$.
\begin{enumerate}
\item Let $\Phi$ be holomorphic on $B(z,r)$. Then there is $M$ with
  $|\Phi^{(k)}(z)|/k!\leq Mr'^{-k}$ for every integer $k\geq0$.
\item Let $\rho\in B(z,r)$ with $\rho\neq z$ and let $c\in\mathbf C$. Let
  $\Phi$ be holomorphic on $B(z,r)\setminus\{\rho\}$ and such that
  $\Phi(s)-c/(s-\rho)$ is holomorphic near $\rho$. Then there is $M$ with
  \[
   \Bigl|\frac{(-1)^k\Phi^{(k)}(z)}{k!}-\frac{c}{(z-\rho)^{k+1}}\Bigr|\leq Mr'^{-k}
  \]
  for every integer $k\geq0$.
\end{enumerate}
\end{lemma}
\begin{proof}
Part (1) is Cauchy's estimate on the circle $|s-z|=r'$, with $M$ the mean of
$|\Phi|$ over that circle. For (2), the function $\Phi(s)-c/(s-\rho)$ extends
to a holomorphic function on $B(z,r)$, to which (1) applies, and the $k$-th
derivative of $c/(s-\rho)$ at $z$ is $(-1)^kk!\,c\,(z-\rho)^{-k-1}$.
\end{proof}
\leanref{Development/AutomorphicArithmetic/DedekindLFunctions/LandauTheorems/LandauCauchy.lean}{Deduction.landau_cauchy}
\leanref{Development/AutomorphicArithmetic/DedekindLFunctions/LandauTheorems/LandauPoleAsym.lean}{Deduction.landau_pole_asym}

\begin{theorem}\label{wl-landau-window}
Let $(g,\Phi,Z)$ be pole data (Definition~\ref{wl-pole-data}). Let
$0<\theta<1$ and $\delta>0$, and suppose that $g$ satisfies the window
hypothesis with exponent $\theta$ and width $\delta$
(Definition~\ref{wl-window-hypothesis}). Then $\operatorname{Re}\rho\leq\theta$
for every $\rho\in Z$.
\end{theorem}
\begin{proof}
Suppose that some $\rho_0\in Z$ has $\operatorname{Re}\rho_0>\theta$. Write
$a_k(s)=(-1)^k\Phi^{(k)}(s)/k!$, and let $\sigma>2$ and $C$ be given by
Lemma~\ref{wl-coefficient-bound}.

Every $\rho\in Z$ has $\operatorname{Re}\rho\leq2<\sigma$. Indeed, let
$|g(u)|\leq C_0e^{2u}$ for $u\geq0$. For $\sigma'>2$ and
$\operatorname{Re}s\geq\sigma'$ the Laplace integral gives
$|\Phi(s)|\leq C_0\int_0^\infty e^{-(\sigma'-2)u}\,du$, so $\Phi$ is bounded on
$\operatorname{Re}s\geq\sigma'$. If $\operatorname{Re}\rho>2$, take
$\sigma'=(2+\operatorname{Re}\rho)/2$: then $(s-\rho)\Phi(s)$ tends to $0$ as
$s\to\rho$, whereas it tends to the nonzero residue of $\Phi$ at $\rho$. So
Lemma~\ref{wl-nearest-pole} applies; write $B(z,r)$ for the open disc of
radius $r$ about $z$.

In case (1) of Lemma~\ref{wl-nearest-pole}, let $D,\gamma,\rho,r$ be as there,
let $c$ be the residue of $\Phi$ at $\rho$, put $z=\sigma+i\gamma$,
$d=|z-\rho|$ and $r'=(d+r)/2$. Then $0<d<r'<r\leq D\leq\sigma-\theta$, where
$d>0$ because $\operatorname{Re}\rho<\sigma$. Every point of $B(z,r)$ or of
$B(\sigma,D)$ has real part greater than $\sigma-D\geq\theta>0$. So $\Phi$ is
holomorphic on $B(z,r)\setminus\{\rho\}$, with $\Phi(s)-c/(s-\rho)$ holomorphic
near $\rho$, and on $B(\sigma,D)$. Lemma~\ref{wl-pole-coefficients} gives $M$
and $M'$ with
\[
 \bigl|a_k(z)-c(z-\rho)^{-k-1}\bigr|\leq Mr'^{-k},\qquad
 |a_k(\sigma)|\leq M'r'^{-k}\qquad(k\geq0).
\]
For $\eta>0$ and $U_0$ let $U\geq U_0$ and $k\geq U$ be given by
Lemma~\ref{wl-coefficient-bound}. Then
\[
 |c|d^{-k-1}\leq|a_k(z)|+Mr'^{-k}
 \leq Ce^{\eta U}(\sigma-\theta)^{-k}-\operatorname{Re}a_k(\sigma)+Mr'^{-k}
 \leq Ce^{\eta U}(\sigma-\theta)^{-k}+(M+M')r'^{-k}.
\]

In case (2), let $\beta,r$ be as there, let $c$ be the residue of $\Phi$ at
$\beta$, put $d=\sigma-\beta$ and $r'=(d+r)/2$, so that
$0<d<r'<r\leq\sigma-\theta$. Since $\beta$ is real, $c$ is real and $c>0$. As
in case (1), Lemma~\ref{wl-pole-coefficients} gives $M$ with
$|a_k(\sigma)-cd^{-k-1}|\leq Mr'^{-k}$ for $k\geq0$. For $\eta>0$ and $U_0$
let $U\geq U_0$ and $k\geq U$ be given by Lemma~\ref{wl-coefficient-bound}.
Taking real parts and using that lemma with $\gamma=0$,
\[
 cd^{-k-1}\leq\operatorname{Re}a_k(\sigma)+Mr'^{-k}
 \leq Ce^{\eta U}(\sigma-\theta)^{-k}-|a_k(\sigma)|+Mr'^{-k}
 \leq Ce^{\eta U}(\sigma-\theta)^{-k}+Mr'^{-k}.
\]
The left side is positive because $c>0$. A pole with negative residue would
make $a_k(\sigma)$ large and negative, which the one-sided bound of
Lemma~\ref{wl-coefficient-bound} allows.

In both cases, since $r'\leq\sigma-\theta$ and $e^{\eta U}\geq1$, there are
$0<d<r'$, $c'>0$ and $A\geq0$ such that for every $\eta>0$ and every $U_0$
there are $U\geq U_0$ and an integer $k\geq U$ with
\[
 c'd^{-k}\leq Ae^{\eta U}r'^{-k};
\]
namely $c'=|c|/d$, and $A=|C|+|M|+|M'|$ in case (1) and $A=|C|+|M|$ in
case (2). Put $\kappa=\log(r'/d)>0$ and $\eta=\kappa/2$. As $k\geq U$,
$Ae^{\eta U}r'^{-k}=Ae^{\kappa U/2-\kappa k}d^{-k}\leq Ae^{-\kappa U/2}d^{-k}$,
which is less than $c'd^{-k}$ once $U$ is large. This contradiction proves the
theorem.
\end{proof}
\leanref{Development/AutomorphicArithmetic/DedekindLFunctions/LandauTheorems/landauWindow.lean}{Deduction.landau_window}
\leanref{Development/AutomorphicArithmetic/DedekindLFunctions/LandauTheorems/LandauPoleReLe.lean}{Deduction.landau_pole_re_le}
\leanref{Development/AutomorphicArithmetic/DedekindLFunctions/LandauTheorems/landauWindow.lean}{Deduction.landau_window_contra}
\leanref{Development/AutomorphicArithmetic/DedekindLFunctions/LandauTheorems/LandauGrowthContra.lean}{Deduction.landau_growth_contra}

\section{Zero-free half-planes}\label{section-wl-zero-free}

\begin{lemma}\label{wl-power-window}
Let $F$ be a character family with weighted Chebyshev sum $G_F$
(Definition~\ref{wl-character-family}). Let $\delta>0$ and
$\theta,C\in\mathbf R$, and suppose that for every $Y_0$ there is $Y\geq Y_0$
with $G_F(y)\leq Cy^\theta$ for all $y\in[Y,Y^{1+\delta}]$. Then the function
$u\mapsto G_F(e^u)$ satisfies the window hypothesis with exponent $\theta$
and width $\delta$ (Definition~\ref{wl-window-hypothesis}).
\end{lemma}
\begin{proof}
Let $\eta>0$ and $U_0$ be given. There is $Y_1\geq1$ with
$|C|\delta\log Y\leq Y^\eta$ for $Y\geq Y_1$. Take $Y\geq\max(e^{U_0},Y_1)$ as
in the hypothesis, put $U=\log Y\geq U_0$ and $R(u)=|C|e^{\theta u}\geq0$. For
$u\in[U,(1+\delta)U]$ we have $e^u\in[Y,Y^{1+\delta}]$, so
$G_F(e^u)\leq Ce^{\theta u}\leq R(u)$. Finally
\[
 \int_U^{(1+\delta)U}R(u)e^{-\theta u}\,du=|C|\delta\log Y\leq Y^\eta=e^{\eta U}.
\]
\end{proof}
\leanref{Development/AutomorphicArithmetic/DedekindLFunctions/LandauTheorems/WindowHypOfRpow.lean}{Deduction.windowHyp_of_rpow}
\leanref{Development/AutomorphicArithmetic/DedekindLFunctions/LandauTheorems/WindowHypUOfWindowHyp.lean}{Deduction.windowHypU_of_windowHyp}

\begin{theorem}\label{wl-zero-free-family}
Let $F=(\chi_i,a_i)_{i\in I}$ be a character family with weighted Chebyshev
sum $G_F$ (Definition~\ref{wl-character-family}). Let $0<\theta<1$,
$\delta>0$ and $C\in\mathbf R$. Suppose that for every $Y_0$ there is
$Y\geq Y_0$ such that
\[
 G_F(y)\leq Cy^\theta\qquad\text{for all }y\in[Y,Y^{1+\delta}].
\]
Then for every $i\in I$ the function $L(s,\chi_i)$ has no zero with
$\operatorname{Re}s>\theta$.
\end{theorem}
\begin{proof}
By Proposition~\ref{wl-pole-data-family} there is $\Phi_F$ such that
$(g_F,\Phi_F,Z_F)$ is pole data, where $g_F(u)=G_F(e^u)$. By
Lemma~\ref{wl-power-window}, $g_F$ satisfies the window hypothesis with
exponent $\theta$ and width $\delta$, so $\operatorname{Re}\rho\leq\theta$ for
every $\rho\in Z_F$ by Theorem~\ref{wl-landau-window}. If $L(s,\chi_i)=0$ with
$\operatorname{Re}s>\theta$, then $\operatorname{Re}s>0$, so $s\in Z_F$ and
$\operatorname{Re}s\leq\theta$, a contradiction.
\end{proof}
\leanref{Development/AutomorphicArithmetic/DedekindLFunctions/LandauTheorems/ZeroFreeOfRpowWindow.lean}{Deduction.zeroFree_of_rpow_window}
\leanref{Development/AutomorphicArithmetic/DedekindLFunctions/LandauTheorems/lemmaL.lean}{Deduction.lemmaL}

\begin{definition}\label{wl-windows}
Let $S\subset\mathbf R$ and $\lambda\in\mathbf R$. The set $S$ has
\emph{windows of exponent $\lambda$} if for every $y\in\mathbf R$ there is
$Y\geq y$ such that every $x$ with $Y\leq x\leq Y^\lambda$ lies in $S$.
\end{definition}
\leanref{Formalism/SharpRestriction/RamificationFrobenius/TotientGaloisClasses/Window.lean}{Formalism.HC.HasWindows}

\begin{theorem}\label{wl-zero-free-level}
Let $K$ be a level with weighted Chebyshev sum $G_K$
(Definition~\ref{wl-weighted-sum}). Let $\lambda>1$ and $0<\theta<1$, and
suppose that for every $y$ there is $Y\geq y$ such that
\[
 G_K(x)\leq x^\theta\qquad\text{for all }x\in[Y,Y^\lambda].
\]
Then for every $\chi\in K$ the function $L(s,\chi)$ has no zero with
$\operatorname{Re}s>\theta$.
\end{theorem}
\begin{proof}
Let $F$ be the character family indexed by $K$ whose member at $\chi\in K$
is $\chi$ itself, as a character modulo its conductor, with weight $a_1=1$
for the trivial character and $a_\chi=w_\chi$ for $\chi\neq1$. It is a
character family: a primitive character is trivial only when its conductor
is $1$, and the weights are positive. By Definition~\ref{wl-weighted-sum},
$G_F=G_K$. Theorem~\ref{wl-zero-free-family} with $\delta=\lambda-1>0$ and
$C=1$ applies, since $[Y,Y^{1+\delta}]=[Y,Y^\lambda]$, and gives the
conclusion.
\end{proof}
\leanref{Development/AutomorphicArithmetic/PrimitiveCosets/Imports.lean}{UpperTheory.Comparison₂.formB}
\leanref{Formalism/SharpRestriction/RamificationFrobenius/TotientGaloisClasses/Window.lean}{Formalism.HC.Level.FormBAt}
\leanref{Development/AutomorphicArithmetic/PrimitiveCosets/LevelFamily.lean}{UpperTheory.Comparison₂.level_family}

\section{Reduction to window bounds along families}\label{section-wl-reduction}

\begin{situation}\label{wl-families}
For each $\varepsilon>0$ let $\mathcal F_\varepsilon$ be a class of families,
the families of \emph{exponent} $\varepsilon$.
Each $F\in\mathcal F_\varepsilon$ has a set of scales $W_F\subset\mathbf R$
and, for each real number $x$, a set of members of $F$ at $x$, each a
definite-side datum (Definition~\ref{th-definite-datum}); such families are
constructed in Chapter~\ref{chapter-assembly}. For $F\in\mathcal F_\varepsilon$, real
numbers $e_D,e_L,K_0,c,e_K,C_L$ and a scale $x$, let a condition
$T_F(e_D,e_L;K_0,c,e_K,C_L;x)$ be given. The family $F$ is \emph{tied on its
scales} with exponents $(e_D,e_L)$ if there are $K_0,c,e_K,C_L$ and $x_0$
such that $T_F(e_D,e_L;K_0,c,e_K,C_L;x)$ holds for every $x\geq x_0$ in
$W_F$. For a level $K$ (Definition~\ref{wl-level}), $F$ \emph{satisfies the
host conditions for $K$} if there are $C$ and $x_0$ such that every member of
$F$ at every $x\geq x_0$ satisfies the host conditions for $K$ at
$(x,\varepsilon,C)$ (Definition~\ref{hf-host-conditions}).
\end{situation}
\leanref{Formalism/SharpRestriction/FamCore.lean}{Formalism.MainNotions.FamCore}
\leanref{Formalism/SharpRestriction/RamificationFrobenius/TotientGaloisClasses/TopDefs.lean}{Formalism.HC.EFamC}
\leanref{Formalism/SharpRestriction/RamificationFrobenius/TotientGaloisClasses/TopWin.lean}{Formalism.HC.ChebyshevInequalityOnFamilyWindows}

\begin{definition}\label{wl-family-conditions}
In Situation~\ref{wl-families}, let $K$ be a level with weighted Chebyshev
sum $G_K$ (Definition~\ref{wl-weighted-sum}), and let $\lambda\in\mathbf R$.
\begin{enumerate}
\item $K$ has the \emph{Chebyshev bound along tied families} if for all real
  $e_D,e_L$, every $\varepsilon>0$ and every $F\in\mathcal F_\varepsilon$ that
  satisfies the host conditions for $K$ and is tied on its scales with
  exponents $(e_D,e_L)$, there is $x_1$ such that
  \[
   G_K(x)\leq x^{1/2+(e_D+e_L+7)\varepsilon}\qquad\text{for every }x\geq x_1\text{ in }W_F.
  \]
\item $K$ has \emph{tied families with windows of exponent $\lambda$} if
  there are real $e_D,e_L$ and $\varepsilon_1>0$ such that for every
  $\varepsilon\in(0,\varepsilon_1)$ some $F\in\mathcal F_\varepsilon$
  satisfies the host conditions for $K$, is tied on its scales with exponents
  $(e_D,e_L)$, and has a set of scales $W_F$ with windows of exponent
  $\lambda$ (Definition~\ref{wl-windows}).
\end{enumerate}
\end{definition}
\leanref{Formalism/SharpRestriction/RamificationFrobenius/TotientGaloisClasses/TopWin.lean}{Formalism.HC.ChebyshevInequalityOnWindows}
\leanref{Formalism/SharpRestriction/RamificationFrobenius/TotientGaloisClasses/TopWin.lean}{Formalism.HC.ChebyshevInequalityAtWindowScale}

\begin{theorem}\label{wl-level-from-windows}
In Situation~\ref{wl-families}, let $K$ be a level with weighted Chebyshev
sum $G_K$ (Definition~\ref{wl-weighted-sum}), and let $\lambda\in\mathbf R$.
Suppose that:
\begin{enumerate}
\item $K$ has the Chebyshev bound along tied families
  (Definition~\ref{wl-family-conditions});
\item $K$ has tied families with windows of exponent $\lambda$;
\item for every $0<\theta<1$ such that the set $\{x:G_K(x)\leq x^\theta\}$
  has windows of exponent $\lambda$ (Definition~\ref{wl-windows}), no
  $L(s,\chi)$ with $\chi\in K$ has a zero with $\operatorname{Re}s>\theta$.
\end{enumerate}
Then $L(s,\chi)\neq0$ for every $\chi\in K$ and every $s$ with
$\operatorname{Re}s>1/2$.
\end{theorem}
\begin{proof}
Let $\chi\in K$ and $\operatorname{Re}s>1/2$. Put
$\tau=\min(\operatorname{Re}s-1/2,1/2)>0$ and $\theta=1/2+\tau/2$, so that
$1/2<\theta<\operatorname{Re}s$ and $\theta<1$. Let $e_D,e_L,\varepsilon_1$
be given by (2), and choose $\varepsilon\in(0,\varepsilon_1)$ with
$(e_D+e_L+7)\varepsilon\leq\tau/2$. Then (2) gives a family
$F\in\mathcal F_\varepsilon$ that satisfies the host conditions for $K$, is
tied on its scales with exponents $(e_D,e_L)$, and has $W_F$ with windows of
exponent $\lambda$. By (1) there is $x_1$ with
\[
 G_K(x)\leq x^{1/2+(e_D+e_L+7)\varepsilon}\leq x^\theta
\]
for every $x\geq\max(x_1,1)$ in $W_F$. A window $[Y,Y^\lambda]\subset W_F$
with $Y\geq\max(x_1,1)$ therefore lies in $\{x:G_K(x)\leq x^\theta\}$, and
such windows exist with $Y$ arbitrarily large. So that set has windows of
exponent $\lambda$, and (3) gives $L(s,\chi)\neq0$.
\end{proof}
\leanref{Formalism/SharpRestriction/RamificationFrobenius/TotientGaloisClasses/TopWin.lean}{Formalism.HC.level_of_tiedWin}
\leanref{Formalism/SharpRestriction/RamificationFrobenius/TotientGaloisClasses/TopWin.lean}{Formalism.HC.exists_small_eps}

\begin{corollary}\label{wl-zero-free-cyclic}
In Situation~\ref{wl-families}, let $\lambda>1$ and let $Q$ be a set of
primitive Dirichlet characters. Suppose that for every $\psi\in Q$ the cyclic
level $K_\psi$, with weighted Chebyshev sum $G_\psi$
(Definition~\ref{wl-cyclic-level}), satisfies:
\begin{enumerate}
\item $K_\psi$ has the Chebyshev bound along tied families
  (Definition~\ref{wl-family-conditions});
\item $K_\psi$ has tied families with windows of exponent $\lambda$;
\item for every $\lambda'>1$ and every $0<\theta<1$ such that the set
  $\{x:G_\psi(x)\leq x^\theta\}$ has windows of exponent $\lambda'$
  (Definition~\ref{wl-windows}), no $L(s,\chi)$ with $\chi\in K_\psi$ has a
  zero with $\operatorname{Re}s>\theta$.
\end{enumerate}
Then $L(s,\psi)\neq0$ for every $\psi\in Q$ and every $s$ with
$\operatorname{Re}s>1/2$.
\end{corollary}
\begin{proof}
Let $\psi\in Q$. Then $\psi\in K_\psi$, and Theorem~\ref{wl-level-from-windows}
applies to $K_\psi$ and $\lambda$, with (3) the case $\lambda'=\lambda$.
\end{proof}
\leanref{Formalism/SharpRestriction/RamificationFrobenius/TotientGaloisClasses/TopWin.lean}{Formalism.HC.zeroFreeOn_of_tiedWin_cyc}
\leanref{Formalism/SharpRestriction/RamificationFrobenius/TotientGaloisClasses/TopWin.lean}{Formalism.HC.zeroFreeOn_of_tiedWin}
\leanref{Formalism/SharpRestriction/RamificationFrobenius/TotientGaloisClasses/ExplicitFormulaErrorOfHostConditions.lean}{Formalism.HC.zeroFreeOn_of_levels}

\begin{remark}\label{wl-supply}
Condition (3) of Theorem~\ref{wl-level-from-windows} holds for every level
when $\lambda>1$, and condition (3) of Corollary~\ref{wl-zero-free-cyclic}
holds for every $\psi$, by Theorem~\ref{wl-zero-free-level}. In
Chapter~\ref{chapter-assembly} the corollary is applied with $Q$ the set of
all primitive Dirichlet characters, the families of definite-side data of
that chapter and their tied condition (Definition~\ref{asm-tied}).
Condition (1) is Theorem~\ref{ib-window-bound-on-scales}. Condition (2)
comes from the existence of tied families for all host requirements
(Definition~\ref{asm-tied-everywhere}) together with hosts with windows of
exponent $\lambda$ (Chapter~\ref{chapter-host-fields}). The result is
Theorem~\ref{asm-zero-free-primitive}, which the later theorems of
Chapter~\ref{chapter-assembly} use with $\lambda=3/2$.
\end{remark}
